% ---------------------------------------------------------------------------
% header.tex — pandoc MD→PDF code block styling
% Pair with: --highlight-style .fcc/pdf/p10k.theme
%
% Approach: tcolorbox owns the block background entirely.
% fvextra patches Highlighting directly to fix the ␣ bug.
% ---------------------------------------------------------------------------

% ---------------------------------------------------------------------------
% Page margins
% ---------------------------------------------------------------------------
\usepackage[top=2cm, bottom=2cm, left=2cm, right=2cm]{geometry}

% ---------------------------------------------------------------------------
% Hyperlinks — controlled via -V colorlinks=true passed to pandoc.
% pandoc's template conditionally sets colorlinks/hidelinks based on that
% variable — no hyperref config needed here.
% ---------------------------------------------------------------------------

% ---------------------------------------------------------------------------
% Images — required for \includegraphics in title page template
% ---------------------------------------------------------------------------
\usepackage{graphicx}
\usepackage{xcolor}
\usepackage{fvextra}
\usepackage{etoolbox}
\usepackage[most]{tcolorbox}

% ---------------------------------------------------------------------------
% p10k colours (xterm-256 exact hex)
% ---------------------------------------------------------------------------
\definecolor{termbg} {HTML}{262626}
\definecolor{termfg} {HTML}{bcbcbc}

% ---------------------------------------------------------------------------
% Fix ␣ bug: patch the Highlighting environment that pandoc defines via
% \DefineVerbatimEnvironment{Highlighting}{Verbatim}{commandchars=...}
% fvextra's \fvset doesn't reach it because the options on DefineVerbatimEnvironment
% take precedence. We redefine it after \begin{document} to inject breaklines.
% ---------------------------------------------------------------------------
\AtBeginDocument{%
  % Redefine Highlighting to add fvextra breaklines + correct foreground colour.
  % Preserve the commandchars option pandoc requires for token coloring.
  \DefineVerbatimEnvironment{Highlighting}{Verbatim}{%
    commandchars=\\\{\},%
    breaklines=true,%
    breakanywhere=false,%
    breaksymbolleft={},%
    fontsize=\small,%
  }%
  % Suppress hyphenation at font level — prevents TeX inserting hyphens on line breaks.
  \BeforeBeginEnvironment{Verbatim}{\hyphenchar\font=-1\relax}%
  \AfterEndEnvironment{Verbatim}{\hyphenchar\font=`\-\relax}%
}

% ---------------------------------------------------------------------------
% Redefine Shaded using tcolorbox — completely independent of framed/snugshade.
% \AtBeginDocument so it fires after pandoc defines Shaded at line ~47.
% ---------------------------------------------------------------------------
\AtBeginDocument{%
  \ifdef{\Shaded}{}{\newenvironment{Shaded}{}{}}%
  \renewenvironment{Shaded}{%
    \begin{tcolorbox}[%
      enhanced,%
      arc=3pt,%
      boxrule=0.5pt,%
      colback=termbg,%
      colframe=termbg!60!black,%
      coltext=termfg,%
      left=6pt, right=6pt, top=4pt, bottom=4pt,%
      before skip=8pt, after skip=8pt,%
      fontupper=\color{termfg}\small,%
    ]%
  }{%
    \end{tcolorbox}%
  }%
}

% ---------------------------------------------------------------------------
% Inline code styling
% Use \scalerel to allow the text to shrink when inside narrow table columns.
% \texttt globally uses \footnotesize so long identifiers fit in pandoc's
% hardcoded longtable column proportions.
% ---------------------------------------------------------------------------
\usepackage{relsize}
\let\oldtexttt\texttt
\renewcommand{\texttt}[1]{%
  \colorbox{termbg}{{\ttfamily\relsize{-1}\color{termfg}#1}}%
}

% ---------------------------------------------------------------------------
% Suppress all hyphenation in prose.
% \pretolerance=10000 makes TeX accept slightly loose lines before trying
% hyphenation. Combined with \hyphenpenalty=10000 it effectively disables
% hyphenation without breaking the layout engine entirely.
% ---------------------------------------------------------------------------
\pretolerance=10000
\hyphenpenalty=10000
\exhyphenpenalty=10000

% ---------------------------------------------------------------------------
% Table row colours — alternating light grey / white body rows, medium grey
% header. Applies to all longtables pandoc generates.
%
% \rowcolors{<start row>}{<odd colour>}{<even colour>}
% Start row 1 so the header row is row 0 (styled separately via \rowcolor).
% \arrayrulecolor keeps the rule lines visible against the light background.
% ---------------------------------------------------------------------------
\usepackage{colortbl}
\definecolor{tableheader}{HTML}{cccccc}
\definecolor{tablerowodd}{HTML}{f5f5f5}
\definecolor{tableroweven}{HTML}{ffffff}

\AtBeginDocument{%
  \rowcolors{2}{tablerowodd}{tableroweven}%
}

% Style the header row — pandoc wraps it in \endfirsthead / \endhead blocks.
% Redefine \toprule and \midrule to inject \rowcolor on the header row.
% \arrayrulecolor darkens the rule lines so they remain visible on light bg.
\arrayrulecolor{gray!50}

% Colour the header row by patching longtable's first-head block.
% \pretocmd injects \rowcolor{tableheader} immediately after \toprule,
% which is where pandoc places the header row.
\usepackage{longtable}
\AtBeginDocument{%
  \pretocmd{\toprule}{\rowcolor{tableheader}}{}{}%
}